% ★★★ 难题

\newcommand{\qSolidRhombusFoldMidpointArcStem}{
已知菱形 $ABCD$ 中，$AB=2$，$\angle BAD=120^\circ$。$E$ 为边 $BC$ 的中点。将 $\triangle ABE$ 沿 $AE$ 翻折成 $\triangle AB_1E$，点 $B_1$ 位于平面 $ABCD$ 上方。连接 $B_1C$ 和 $B_1D$，$F$ 为 $B_1D$ 的中点。
\begin{enumerate}
    \item 在翻折过程中，$AE$ 与 $B_1C$ 的夹角大小是否改变？若不改变，求出其大小；若改变，说明理由；
    \item 求点 $F$ 的轨迹长度。
\end{enumerate}
}

\newcommand{\qSolidRhombusFoldMidpointArcOptions}{
}

\newcommand{\qSolidRhombusFoldMidpointArcFigure}[1][1]{
\begin{tikzpicture}[scale=#1, x={(-0.5cm,-0.4cm)}, y={(1cm,0cm)}, z={(0cm,1cm)}]
    \coordinate (A) at (0, 0, 0);
    \coordinate (D) at (0, 2, 0);
    \coordinate (B) at (1.732, -1, 0);
    \coordinate (C) at (1.732, 1, 0);
    \coordinate (E) at (1.732, 0, 0);
    
    \coordinate (B1) at (1.732, -0.5, 0.866);
    \coordinate (F) at (0.866, 0.75, 0.433);
    
    \draw[dashed, gray] (B) -- (C) -- (D) -- (A) -- cycle;
    \draw[dashed, gray] (B) -- (E);
    
    \draw (A) -- (D) -- (C) -- (E);
    \draw[dashed] (A) -- (E);
    \draw (A) -- (B1) -- (E);
    \draw (B1) -- (C);
    \draw[dashed] (B1) -- (D);
    
    \node[left] at (A) {$A$};
    \node[right] at (C) {$C$};
    \node[right] at (D) {$D$};
    \node[below, gray] at (B) {$B$};
    \node[below] at (E) {$E$};
    \node[above] at (B1) {$B_1$};
    \node[above right] at (F) {$F$};
    \fill (F) circle (1.5pt);
\end{tikzpicture}
}

\newcommand{\qSolidRhombusFoldMidpointArcAnswer}{
(1) 不改变，$\frac{\pi}{2}$；(2) $\frac{\pi}{2}$
}

\newcommand{\qSolidRhombusFoldMidpointArcSolution}{
\begin{enumerate}
    \item \textbf{不改变，夹角为 $\frac{\pi}{2}$。} \\
    由 $AB=2$，$\angle BAD=120^\circ$，可得 $\angle ABE = 60^\circ$。
    因为 $E$ 为 $BC$ 的中点，且 $BC=2$，所以 $BE=1$。
    在 $\triangle ABE$ 中，由余弦定理求得 $AE=\sqrt{3}$，且 $AE^2+BE^2 = 3+1 = 4 = AB^2$，所以 $AE \perp BE$。
    翻折后，$AE \perp B_1E$，又 $AE \perp CE$，
    且 $B_1E \cap CE = E$，$B_1E, CE \subset \text{平面 } B_1EC$，所以 $AE \perp \text{平面 } B_1EC$。
    又因为 $B_1C \subset \text{平面 } B_1EC$，所以 $AE \perp B_1C$。
    即翻折过程中，$AE$ 与 $B_1C$ 的夹角始终为 $\frac{\pi}{2}$。
    
    \item \textbf{轨迹长度为 $\frac{\pi}{2}$。} \\
    点 $B_1$ 的运动轨迹是以 $E$ 为圆心，在垂直于 $AE$ 的平面内，半径为 $B_1E=1$ 的半圆。
    因为 $F$ 为 $B_1D$ 的中点，由中位线定理或位似变换可知，点 $F$ 的轨迹是点 $B_1$ 轨迹以 $D$ 为位似中心，相似比为 $\frac{1}{2}$ 的图形。
    因此，$F$ 的轨迹也是半圆，其半径为 $\frac{1}{2}$。
    所以，点 $F$ 的轨迹长度为 $\frac{1}{2} \times \pi \times 1 = \frac{\pi}{2}$。
\end{enumerate}
}

\newcommand{\qSolidRhombusFoldMidpointArcDiagnosis}{
\textbf{学生常见问题：}
\begin{itemize}
    \item \textbf{空间想象障碍}：无法发现 $AE \perp BC$，导致对翻折后的垂直关系不敏感，建立不出以 $AE$ 为轴的旋转模型。
    \item \textbf{轨迹缩放不清}：知道 $B_1$ 的轨迹是圆弧，但在求 $F$ 的轨迹时，忽略了位似变换或中点平移，无法直接得出轨迹半径为 $\frac{1}{2}$，从而采用建系硬算，导致时间浪费。
\end{itemize}
}

\newcommand{\qSolidRhombusFoldMidpointArcTeachingNotes}{
\textbf{课堂处理与追问：}
\begin{itemize}
    \item \textbf{模型建立}：引导学生先算平面几何中的特殊角和垂直关系，看到 $1, \sqrt{3}, 2$ 必须想到直角，这是折叠问题的基石（找不变的线面垂直）。
    \item \textbf{追问 1}：点 $B_1$ 形成的半圆，其所在平面与 $AE$ 的关系是什么？（垂直，因为 $AE \perp$ 平面 $B_1EC$）。
    \item \textbf{追问 2}：除了用位似，如果用空间向量如何描述 $F$ 的坐标？（借此复习圆在空间中的参数方程表示法）。
\end{itemize}
}
